\item \textbf{Problema de repaso.} Dos balas de plomo, una de masa 12.0 g con una rapidez de 300 m/s se mueve hacia la derecha, y otra de masa 8.00 g con una rapidez de 400 m/s se desplaza hacia la izquierda, colisionan de frente y permanecen juntas. Ambas balas están originalmente a $30.0^\circ\text{C}$. Suponga que el cambio de energía cinética del sistema aparece totalmente como energía interna aumentada. Nos interesa determinar la temperatura y fase de las balas después de la colisión. (a) ¿Cuáles son dos modelos de análisis apropiados para el sistema de dos balas, durante el intervalo de tiempo desde antes hasta después de la colisión? (b) De uno de esos modelos, ¿cuál es la rapidez de las balas combinadas después de la colisión? (c) ¿Cuánta de la energía cinética inicial se ha transformado a energía interna en el sistema después de la colisión? (d) ¿Se funde todo el plomo debido a la colisión? (e) ¿Cuál es la temperatura de las balas combinadas después de la colisión? (f) ¿Cuál es la fase de las balas combinadas después de la colisión?
